% ARCHIVED DRAFT: independently audited by root and incorporated 2026-11-19 18:10 UTC.
% Main part_B.tex is authoritative.
\begin{saresult}{B.081}{Proved under stated hypotheses}{Self-injective finite approximations and the unbounded nilpotent test ring}
Every injective ring epimorphism $R\to C$ is an isomorphism in either
of the following situations:
\begin{enumerate}
\item $R$ is a directed union of unital self-injective subrings, where
self-injective means injective as a module over itself.
\item $R$ is local with nil maximal ideal $\mathfrak m$, and it has
an ideal $I$ such that $\mathfrak m^N I=0$ for some finite $N$, while
every injective ring epimorphism out of $R/I$ is an isomorphism.
\end{enumerate}
For any such base and any unital inclusion $R\subset S$, the polynomial
co-operations exist uniquely exactly when
\[
 (R:_S D_d(R))=R\quad\text{for all }d\ge1,
 \quad\text{equivalently}\quad \Omega(R,S)=R[X].          \tag{B81.1}
\]
Every positive case has all multivariable tensor isomorphisms.
In particular this completely decides all coefficient extensions of
the previously specified test base
\[
 R=\mathbb F_2[z,x_1,x_2,\ldots]/
       (z^2,\ zx_i\ (i\ge1),\ x_i^{2^{i+1}}\ (i\ge1)).   \tag{B81.2}
\]
This base is local of dimension zero and non-Noetherian; its maximal
ideal is not T-nilpotent and its elementwise nilpotence has no uniform
bound. Thus the conclusion does not follow from those earlier
bounded-nilpotence cases.
\end{saresult}
\begin{proof}
For the first assertion write $R=\bigcup_\alpha R_\alpha$ as a directed
union, and suppose $R\subset C$ is epic. For $c\in C$ the equality
$c\otimes1=1\otimes c$ holds in $C\otimes_R C$. The balancing
relations defining this tensor product are the union of those for
the $R_\alpha$: every individual equality is a finite sum of
balancing relations and therefore uses coefficients from one
$R_\alpha$. Consequently
\[
          c\otimes1=1\otimes c
                   \quad\text{in }C\otimes_{R_\alpha}C
                                                        \tag{B81.3}
\]
for some $\alpha$. Since $R_\alpha$ is self-injective, its inclusion
as an $R_\alpha$-submodule of $C$ has an $R_\alpha$-linear retraction
$\rho:C\to R_\alpha$. In particular $\rho(1)=1$.
Applying $\rho\otimes1$ to (B81.3) gives $\rho(c)=c$ in $C$.
Hence $c\in R_\alpha\subset R$, proving $C=R$. This argument uses
neither flatness of the transition maps nor a module-finiteness
assumption on $C$.

For the second assertion, let again $R\subset C$ be epic. The ideal
$\mathfrak mC$ is nil: each element is a finite sum of multiples of
nilpotent elements, and the ideal generated by any finite collection
of such elements is nilpotent in a commutative ring. Hence
$C/\mathfrak mC$ is nonzero. Base change gives an epimorphism from
the field $R/\mathfrak m$ to this nonzero algebra, so it is an
isomorphism. For completeness a field extension algebra $B$ with
$B\otimes_kB\simeq B$ must equal $k$: for any $b\notin k$ a
$k$-linear functional vanishing on $1$ and nonzero on $b$ contradicts
$b\otimes1=1\otimes b$. We conclude $C=R+\mathfrak mC$.
Multiplying by $I$ and iterating this equality gives
\[
                 IC=I+\mathfrak mIC
                    =I+\mathfrak m^NIC=I.               \tag{B81.4}
\]
Thus $I$ is an ideal of $C$ with precisely its original underlying
subset of $R$. The induced map $R/I\to C/I$ is therefore an
injective ring epimorphism. By hypothesis it is an isomorphism, and
then $C=R+I=R$.

For any base with this epimorphism property, B.064 applies to the
explicit effective coefficient ring $C=\bigcup_d(R:_S D_d)$.
Tensor injectivity forces $R\to C$ epic, hence $C=R$. Conversely
$C=R$ gives ordinary polynomials in all arities. This proves (B81.1)
and all positive tensor assertions.

Finally put
\[
 A_n=\mathbb F_2[x_1,\ldots,x_n]/
                (x_i^{2^{i+1}}:1\le i\le n),
 \qquad A=\bigcup_{n\ge1}A_n.                            \tag{B81.5}
\]
The transition maps are the evident injections on the monomial
bases. Every $A_n$ is self-injective. Indeed the functional taking
the coefficient of the top monomial
$\prod_{i=1}^n x_i^{2^{i+1}-1}$ gives a nondegenerate bilinear form
$(a,b)\mapsto\ell(ab)$: a monomial pairs with its complementary
monomial, and the resulting matrix is a permutation matrix.
It identifies $A_n$ with $\operatorname{Hom}_{\mathbb F_2}
(A_n,\mathbb F_2)$ as an $A_n$-module. This dual module is injective,
since for every $A_n$-module $M$ there is a natural identification
\[
 \operatorname{Hom}_{A_n}
       (M,\operatorname{Hom}_{\mathbb F_2}(A_n,\mathbb F_2))
       \simeq\operatorname{Hom}_{\mathbb F_2}(M,\mathbb F_2),
\]
and vector-space duality is exact. The first part of the theorem
therefore proves that all injective epimorphisms out of $A$ are
isomorphisms.

The ring (B81.2) is $A\oplus\mathbb F_2z$, where $z^2=0$ and the
augmentation ideal of $A$ annihilates $z$. Its maximal ideal
$\mathfrak m=(z,x_1,x_2,\ldots)$ is nil, since any element uses only
finitely many variables, and the corresponding finite algebra has
nilpotent maximal ideal. Thus $R$ is local of dimension zero.
The ideal $I=\mathbb F_2z$ satisfies $\mathfrak mI=0$ and $R/I=A$,
so the second part applies. The classes of all $x_i$ are independent
in $\mathfrak m/\mathfrak m^2$, proving non-Noetherianity. Each product
$x_1\cdots x_n$ is a nonzero basis monomial, proving failure of
T-nilpotence. The exact nilpotence orders $2^{i+1}$ of the $x_i$ are
unbounded, which checks the remaining assertions about the test ring.
\end{proof}
